\BourbakiExerciseGroup

\begin{enumerate}
\def\labelenumi{\arabic{enumi})}
\tightlist
\item
  设 \(K_{0}\) 是特征为 p \textgreater{} 0 的域，\(\mathbf{K} = \mathbf{K}_{0}(\mathbf{X}, \mathbf{Y})\) 是 \(K_{0}\) 上两个未定元的有理分式域，U 与 V 是两个未定元，\(\Omega\) 是 \(\mathbf{K}(\mathbf{U}, \mathbf{V})\) 的一个代数闭扩张，\(\mathbf{E} = \mathbf{K}(\mathbf{U}, u)\) ，其中 \(u^{p} = X + YU^{p}\) ，而 \(F = K(V, v)\) ，其中 \(v^{p} = X + YV^{p}\) ；E 与 F 在 K 上不是可分的（V，第 178 页，习题 6，c））。
\end{enumerate}

\begin{enumerate}
\def\labelenumi{\alph{enumi})}
\item
  证明 \(\mathbf{K}\) 在 \(\mathbf{E}\) 中与在 \(\mathbf{F}\) 中都是相对代数封闭的（若 \(x \in \mathbf{E}\) 在 \(\mathbf{K}\) 上是代数的，证明必有 \(x^p \in \mathbf{K}\) ；若 \(x \notin \mathbf{K}\) ，则 \(\mathbf{E} = \mathbf{K}(\mathbf{U}, x)\) ，由此推出 \(\mathbf{X}^{1/p}\) 与 \(\mathbf{Y}^{1/p}\) 将属于 \(\mathbf{K}(x)\) ）。
\item
  证明 E 与 F 在 K 上线性不相交，但 K 在 \(\mathbf{K}(\mathbf{E} \cup \mathbf{F})\) 中不是相对代数封闭的。（证明 \(X^{1/p} \in \mathbf{K}(\mathbf{E} \cup \mathbf{F})\) ；由此推出 v 不能属于 \(\mathbf{E}(\mathbf{V})\) ，并由此得出 E 与 F 在 K 上线性不相交。）
\end{enumerate}

\begin{enumerate}
\def\labelenumi{\arabic{enumi})}
\setcounter{enumi}{1}
\tightlist
\item
  设 \(\mathbf{K}\) 是特征为 \(p > 0\) 的域，\(\Omega\) 是 \(\mathbf{K}\) 的一个扩张，而 \(\mathbf{E}\) 、\(\mathbf{F}\) 是 \(\Omega\) 的两个在 \(\mathbf{K}\) 上代数不相交的子扩张。
\end{enumerate}

\begin{enumerate}
\def\labelenumi{\alph{enumi})}
\tightlist
\item
  设 \(F_{s}\) 是 K 在 F 中的相对可分闭包。证明 \(E(F_{s})\) 是 E 在 \(E(F)\) 中的相对可分闭包。（归结为 \(F_{s}=K\) 的情形；设 B 是 E 在 K 上的一个超越基，而 \(K'\) 是 K 在 F 中的相对代数闭包，它在 F 上是根式的；则（V，第 137 页）\(K'(B)\) 是 \(K(B)\) 在 \(F(B)\) 中的相对代数闭包，且在 \(K(B)\) 上是 p-根式的。若 \(x\in E(F)\) 在 E 上是代数的，则存在整数 \(r\geqslant0\) 使得 \(x^{p'}\in M\) ，其中 M 是 \(F(B)\) 的一个有限次可分扩张，具有一组基 \((u_{i})_{1\leqslant i\leqslant n}\) （在 \(\mathrm{F(B)}\) 上，由 E 的元素组成）。若 \(y = \sum_{i}b_{i}u_{i}\) ，其中 \(b_{i}\in \mathrm{F(B)}\) ，证明
\end{enumerate}

通过考虑 M 到 E(F) 的一个代数闭包中的 F(B)-同构，证明 \(b_{i}\) 属于 K'(B)，并由此推出 y 在 E 上是 p-根式的。）

\begin{enumerate}
\def\labelenumi{\alph{enumi})}
\setcounter{enumi}{1}
\tightlist
\item
  假设 E 与 F 在 K 上线性不相交，K 在 F 中相对代数封闭，且 E 是 K 的可分扩张。证明 E 在 E(F) 中是相对代数封闭的。（借助 V，第 138 页，归结为 E 是 K 的可分代数扩张的情形。由 a) 推出 E 在 E(F) 中的相对代数闭包在 E 上是 p-根式的，然后利用如下事实：若 \((a_{\lambda})\) 是 E 在 K 上的一组基，则诸 \(a_{\lambda}^{p'}\) 也构成 E 在 K 上的一组基。）
\end{enumerate}

\begin{enumerate}
\def\labelenumi{\arabic{enumi})}
\setcounter{enumi}{2}
\tightlist
\item
  设 E 是一个域，G 是 E 的一个自同构群，而 K 是 G 的不变域。a) 证明：若 G 中不存在异于 G 本身的有限指数正规子群，则 E 是 K 的正则扩张。
\end{enumerate}

\begin{enumerate}
\def\labelenumi{\alph{enumi})}
\setcounter{enumi}{1}
\tightlist
\item
  由 a) 推出：若 G 同构于无限单群的一个有限直积，则 E 是 K 的正则扩张。
\end{enumerate}

\begin{enumerate}
\def\labelenumi{\arabic{enumi})}
\setcounter{enumi}{3}
\tightlist
\item
  设 K 为一个域，E、F 为 K 的两个扩张。
\end{enumerate}

\begin{enumerate}
\def\labelenumi{\alph{enumi})}
\item
  假设 K 在 F 中相对代数封闭。证明：若 \((L, u, v)\) 与 \((L', u', v')\) 是 E 与 F 的两个复合扩张，使得 \(u(E)\) 与 \(v(F)\) （相应地 \(u'(E)\) 与 \(v'(F)\) ）在 K 上代数不相交，则这两个复合扩张是同构的。（首先考虑 E 是 K 的纯超越扩张的情形；然后利用习题 2，a)，以及 V，第 154 页，习题 6，b)）
\item
  设 \(E_{s}\) 、\(F_{s}\) 分别是 K 在 E 与 F 中的相对可分闭包。假设 \(E_{s}\) 与 \(F_{s}\) 的所有复合扩张都是同构的。证明：若 \((L, u, v)\) 与 \((L', u', v')\) 是 E 与 F 的两个复合扩张，使得 \(u(E)\) 与 \(v(F)\) （相应地 \(u'(E)\) 与 \(v'(F)\) ）在 K 上代数不相交，则这两个复合扩张是同构的（利用 a) 以及 V，第 173 页，习题 11）。
\end{enumerate}

\begin{enumerate}
\def\labelenumi{\arabic{enumi})}
\setcounter{enumi}{4}
\tightlist
\item
  设 K 是一个域，p 是它的特征指数。对每个与 p 互素的整数 n \textgreater{} 0 ，用 \(\varphi_{\mathrm{K}}(n)\) 表示 \(\mathrm{R}_{n}(\mathrm{K}) = \mathrm{K}(\mu_{n})\) 在 K 上的次数，其中 \(\mu_{n}\) 表示 K 的一个代数闭包中 n 次单位根所成的群。例如 \(\varphi_{\mathrm{Q}}(n) = \varphi(n)\) （V，第 84 页，定理 2）；若 K 是含 q 个元素的有限域，则 \(\varphi_{\mathrm{K}}(n)\) 是 q 的剩余类在 \((\mathbf{Z}/n\mathbf{Z})^{*}\) 中的阶（V，第 97 页，推论）；若 K 是代数闭的，则对所有 n 有 \(\varphi_{\mathrm{K}}(n) = 1\) 。
\end{enumerate}

\begin{enumerate}
\def\labelenumi{\alph{enumi})}
\item
  设 E 是素域 \(K_{0}\) 在 K 中的相对代数闭包。证明 \(\varphi_{\mathrm{K}}(n) = \varphi_{\mathrm{E}}(n)\) 。
\item
  假设 K 是 \(K_{0}\) 的有限生成扩张。设 m 是整数 \textgreater0 ；则存在整数 N \textgreater0 ，它被每个满足 \(\varphi_{\mathrm{K}}(n) \leqslant m\) 的 n 整除。
\item
  设 \(s\) 是 \(\mathrm{GL}_m(\mathbf{K})\) 的一个有限阶元素。证明 \(s^N\) 的阶是 \(p\) 的一个幂。（我们有 \(s = tu = ut\) ，其中 \(u\) 的阶是 \(p\) 的幂，而 \(t\) 的阶 \(n\) 与 \(p\) 互素。元素 \(t\) 是可分多项式 \(\mathbf{X}^n - 1\) 的根，同时也是其特征多项式的根；由此推出子代数 \(\mathbf{A}\) —— \(\mathbf{M}_m(\mathbf{K})\) 的、由 \(t\) 生成的 —— 是平展的，次数 \(\leqslant m\) ，同构于一些扩张 \(\mathbf{R}_d(\mathbf{K})\) 的积，其中 \(d\) 遍历 \(n\) 的某些因子，包括 \(n\) 本身。因此 \(\varphi_{\mathbf{K}}(n) \leqslant m\) ；利用 \(b\) ）。）
\end{enumerate}
% semantic-fragments: legacy-input-seam
\relax{}
